\scanpage{101}
\begin{center}- \mspage{171} -\end{center}
\margin{IX}
engendre $R$ au voisinage de $s$, d'où il résulte que $F(R)\subset R'$ au voisinage de $s$, ce qui achève la démonstration.

\uline{Corollaire} 13.2.2. \uline{Si} $A$ \uline{satisfait les conditions} a), b), c) \uline{de} 13.2.2, \uline{il est déterminé à isomorphisme unique près au voisinage de} $s$ \uline{par la connaissance du} $\widehat S$-\uline{groupe} $\widehat A$.

13.3. \uline{L'extension de Raynaud analytique}. Nous sommes maintenant en mesure de paraphraser, en Géométrie Analytique, la construction de l'extension de Raynaud donnée dans 7. Supposons donné $A$ comme dans 13.1, avec $A=A_0$, et considérons $R_0=R_s$, qui est un $\mathbb Z$-module libre de type fini. Quitte à restreindre $S$ à un voisinage ouvert assez petit de $s$, interprétant $R_0$ comme le groupe des germes de sections de $R$ au voisinage de $s$, on trouve $\overline{\text{un}}$ homomorphisme
\[
(13.3.1)\qquad R_{0S}\longrightarrow R\subset E
\]
où $R_{0S}$ désigne le groupe analytique constant sur $S$ de valeur $R_0$. Choisissant le voisinage de $s$ assez petit pour que toute section de $\underline{\mathrm O}_S$ dont le germe en $s$ est nul soit nul, on voit que l'homomorphisme précédent est injectif. Son image est donc un \uline{sous}-groupe analytique constant de $R$, qui est évidemment le plus grand sous-groupe constant de $R$, et qui s'appelle la partie \uline{fixe} de $R$, et est notée $R^f$ :
\[
(13.3.2)\qquad R_{0S}\simeq R^f\unsure{\subset|\hookrightarrow}R,\ \text{induit un isomorphisme }R^f_s\simeq R_s.
\]
On pourra alors considérer le groupe quotient
\[
(13.3.3)\qquad A^{\natural}=E/R^f
\]
\margin{\unsure{489}}
\margin{489}
\margin{100}
